% ============================================================
%  Chapter 2 — Methodology
%  Ported from writeup/methodology/methodology_overview.tex §§3–10
%  (June 2026). Data-facing material (state derivation, macro data
%  sources, exploratory evidence) lives in Chapter 3 and is
%  cross-referenced, not duplicated.
% ============================================================

\chapter{Methodology}
\label{chap:methodology}

\section{The transition model}
\label{sec:meth-model}

The supervised object is the one-month transition distribution
\eqref{eq:object} over the seven states defined in
Section~\ref{sec:data-outcome}.  In this dataset the terminal states
(\emph{Prepaid}, \emph{Foreclosure}, \emph{REO}) are absorbing ---
they are observed only in a loan's final month --- so the origin set
comprises the four live states and the transition matrix is
$4 \times 7$ rather than $7 \times 7$ (a documented difference from
\citealp{sirignano2021}, harmless for the model comparison since every
competing model conditions on the same origin set).  A single
multinomial model is fitted with the current state as an input
feature, rather than four separate per-origin models.  Some
transitions are structurally impossible --- delinquency deepens at
most one bucket per month, so \emph{Current} cannot reach 60\,DPD
directly --- and rather than hard-coding a mask, the models are left
to learn these near-zeros; the evaluation suite asserts that the
predicted mass on impossible cells is negligible, a useful internal
consistency check.

\section{Estimation at scale: sharding and minibatch randomisation}
\label{sec:meth-shards}

Stochastic gradient descent assumes minibatches are approximately
random draws from the training distribution.  The panel, however, is
stored grouped by acquisition vintage: read sequentially, it yields
batches composed entirely of, say, 2007-vintage loans, whose risk
profile is wildly atypical, and the resulting gradient bias changes
within an epoch.  The textbook remedy --- globally shuffling the
dataset each epoch --- is infeasible for $\sim$$10^9$ rows on external
storage.

The pipeline instead assigns every row a \emph{shard},
\begin{equation}
\label{eq:shard}
\mathrm{shard}(i) \;=\; h\!\left(\mathrm{LoanID}_i\right) \bmod N,
\qquad N=256,
\end{equation}
where $h$ is a fixed-seed 64-bit hash.  Three properties follow, each
load-bearing.  First, a good hash behaves like an independent uniform
draw per loan identifier, so each shard is an unbiased $\tfrac1N$
sample of \emph{loans}, mixing all vintages and geographies (verified
empirically: rows and exact distinct loans per shard are uniform to
coefficients of variation of $0.26\%$ and $0.21\%$, and per-shard
vintage composition matches the panel marginal).  Second, because the
hash is keyed on the loan --- not the row --- every monthly
observation of a loan lands in the same shard, so shard-based subsets
never split a loan between training and evaluation.  Third, the
assignment is a pure function of the loan identifier: rebuilding the
panel reproduces every shard exactly, with no random state to
persist.

\begin{figure}[htbp]
\centering
\begin{tikzpicture}[node distance=4mm and 6mm,
  box/.style={draw, rounded corners=1pt, align=center, font=\small, inner sep=4pt},
  arr/.style={-{Stealth[length=2mm]}}]
\node[box] (a) {epoch start:\\ shuffle shard\\ \emph{order}};
\node[box, right=of a] (b) {load one shard\\ (RAM-sized,\\ all vintages mixed)};
\node[box, right=of b] (c) {shuffle rows\\ \emph{within}\\ the shard};
\node[box, right=of c] (d) {yield\\ minibatches,\\ SGD steps};
\node[box, right=of d] (e) {next shard\\ until epoch\\ ends};
\draw[arr] (a) -- (b);
\draw[arr] (b) -- (c);
\draw[arr] (c) -- (d);
\draw[arr] (d) -- (e);
\end{tikzpicture}
\caption[The two-level shuffle over loan-keyed shards]{\textit{The
two-level shuffle.}  Random shard order across the epoch, full
shuffle within each RAM-sized shard.  Every minibatch mixes loans
from all vintages at streaming cost; a global shuffle of $10^9$ rows
is never performed.}
\label{fig:shardloop}
\end{figure}

Training then follows the two-level shuffle of
Figure~\ref{fig:shardloop}.  It is not exactly a uniform draw --- rows
in a batch share a shard --- but because each shard is already an
unbiased cross-section of loans, the residual within-shard
correlation is negligible for gradient quality.

\section{Covariates and feature construction}
\label{sec:meth-features}

Features comprise the loan-level fields and external macroeconomic
covariates described in Sections~\ref{sec:data-variables}
and~\ref{subsec:data-external}.  Continuous covariates are
standardised using moments estimated on each training window only
(Section~\ref{sec:meth-design}); heavily skewed balances enter in
logarithms.  High-cardinality categorical fields (property state,
metropolitan area, three-digit zip code) are mapped to learned dense
embedding vectors rather than one-hot columns, with categories unseen
in training routed to a reserved index.  Missingness indicators
accompany sentinel-scrubbed fields, since the fact of missingness is
itself predictive in credit data.

Two engineered covariates carry most of the economic content:
\begin{align}
\label{eq:incentive}
\text{(refinancing incentive)} \qquad
\mathrm{incentive}_{i,t} &= r^{\mathrm{cur}}_{i,t} - m(t),\\[2pt]
\label{eq:mtmltv}
\text{(mark-to-market LTV)} \qquad
\mathrm{ltv}^{\mathrm{mtm}}_{i,t} &=
\frac{\mathrm{UPB}_{i,t}}{\mathrm{UPB}_{i,0}}\;\cdot\;
\mathrm{LTV}_{i,0}\;\cdot\;
\frac{H_{s(i)}\!\left(t_0(i)\right)}{H_{s(i)}(t)},
\end{align}
where $m(t)$ is the prevailing 30-year mortgage rate,
$\mathrm{UPB}$ denotes the unpaid principal balance (so
$\mathrm{UPB}_t/\mathrm{UPB}_0$ is the fraction of the loan still
owed), and $H_s(\cdot)$ is the state-$s$ house-price index at the
indicated month.  Equation~\eqref{eq:incentive} measures how far the
borrower's refinancing option is in the money --- the dominant
prepayment driver \citep{schwartz1989} --- and
\eqref{eq:mtmltv} approximates the borrower's current leverage, the
key default driver \citep{deng2000}.  Raw calendar time is
deliberately excluded as a feature: test years are by construction
unseen, so a model leaning on ``year'' would extrapolate arbitrarily;
all time variation must flow through observable, recurring channels
(the macro block, seasonality encoded as month-of-year, and loan
age).

\section{Rolling out-of-sample design}
\label{sec:meth-design}

Credit models are used prospectively, so the evaluation is temporal:
train on the past, predict the future.  One subtlety deserves
emphasis because it is a one-character bug with real leakage
consequences: the label of the example built at feature month $t$ is
the state at $t{+}1$, so ``training on data up to month $T$'' must
mask on the \emph{label} month --- include the example only if
$t < T$ strictly, so that its label month $t{+}1 \le T$.

The primary design is a rolling, expanding-window backtest with one
window per test year $k = 2015,\dots,2025$:
\[
\underbrace{\text{train: labels} \le \mathrm{Dec}(k{-}2)}_{\text{expanding, starts 2000}}
\qquad
\underbrace{\text{validation: labels in year } k{-}1}_{\text{early stopping only}}
\qquad
\underbrace{\text{test: labels in year } k}_{\text{never touched in fitting}}
\]
Three rules keep this honest and affordable.  First,
\emph{tune once, freeze, roll}: architecture and regularisation
hyperparameters are selected once, on the $k{=}2015$ tuning window,
then frozen across all windows; per-window validation years serve
only early stopping.  Re-tuning every window would multiply compute
by the grid size and constitute a subtle multiple-testing problem.
Second, \emph{window-local preprocessing}: feature scalers, embedding
vocabularies, and any statistic estimated from data are re-fitted
within each window's training slice; standardising with full-sample
moments would leak future means and variances into the past.  Third,
\emph{pooled and per-year reporting}: each test row in 2015--2025 is
predicted exactly once, by a model that never saw it, and metrics are
reported per test year and pooled over all years.  Regime shifts ---
COVID in 2020--21, the rate shock of 2022--23 --- thus appear as
legible test-year effects rather than hinging on a single arbitrary
cutoff.

\paragraph{Class thinning with importance weights.}
\emph{Current}-to-\emph{Current} transitions are $96.3\%$ of
loan-months (Table~\ref{tab:stateshares}) and individually nearly
uninformative.  The training export therefore retains every row whose
origin or destination is not \emph{Current}, but keeps
\emph{Current}$\to$\emph{Current} rows only with probability $p$,
attaching weight $w_i = 1/p$ to kept rows ($w_i = 1$ otherwise).  The
training objective is the weighted negative log-likelihood
\begin{equation}
\label{eq:weightednll}
\widehat{\mathcal{L}}(\theta) \;=\;
\frac{1}{\sum_i w_i}\sum_{i \in \text{kept}} w_i\,
\bigl[-\log p_\theta\!\left(y_i \mid x_i\right)\bigr],
\end{equation}
a Horvitz--Thompson estimator of the full-sample loss
\citep{ht1952}: each row is kept with known probability and weighted
by its inverse, so the fitted model still estimates the \emph{true}
conditional probabilities --- unlike naive downsampling, which would
silently inflate every rare-event probability.  Validation and test
slices are never thinned, and a QA step compares ensemble-average
predicted base rates against unthinned panel frequencies.
Evaluation uses a fixed, loan-disjoint shard block ($\sim$20\% of
loans, the same loans in every window and for every model), so that
likelihood comparisons are exact rather than sampled apart.

\section{The model hierarchy}
\label{sec:meth-models}

The benchmarks are ordered by how flexibly they let
\eqref{eq:object} depend on covariates; each strictly nests or
dominates the previous, which is what makes the comparison
interpretable.

\subsection{Empirical transition matrix}
\label{subsec:meth-empirical}

The simplest estimator ignores covariates entirely: conditional only
on the origin state $u$, the destination distribution is estimated by
smoothed relative frequencies in the window's training slice,
\begin{equation}
\label{eq:empirical}
\widehat{P}_{uv} \;=\; \frac{N_{uv} + \alpha}{N_{u\cdot} +
\alpha\,|\mathcal{S}|}, \qquad \alpha = \tfrac12,
\end{equation}
where $N_{uv}$ counts $u\to v$ transitions.  The Laplace term
$\alpha$ guarantees no test transition receives probability zero (a
single unsmoothed empty cell would make the test log-likelihood
infinite).  This time-homogeneous Markov chain is the floor of the
comparison.  A ``bucketed'' variant conditioning additionally on
credit-score and incentive terciles shows how far cheap,
nonparametric conditioning goes before any model fitting.

\subsection{Multinomial logistic regression as a zero-layer network}
\label{subsec:meth-logit}

The workhorse of the empirical mortgage literature posits log-linear
odds:
\begin{equation}
\label{eq:logit}
p_\theta(j \mid x) \;=\; \operatorname{softmax}_j\!\left(Wx + b\right)
\;=\; \frac{\exp\!\left(w_j^\top x + b_j\right)}
{\sum_{j'} \exp\!\left(w_{j'}^\top x + b_{j'}\right)}.
\end{equation}
Crucially, \eqref{eq:logit} is exactly a neural network with zero
hidden layers, and it is implemented as such --- same data loader,
same loss \eqref{eq:weightednll}, same evaluation code as the deep
models, with a one-time numerical cross-check against a standard
library implementation.  This is an experimental-control choice: when
the deep network beats the logit, the difference is attributable to
architecture alone, not to differences in preprocessing, optimisation
or data.  An ``augmented'' logit with hand-crafted nonlinear terms
(squared loan age, splined incentive) quantifies how much of the gap
manual feature engineering can close.

\subsection{Deep neural networks}
\label{subsec:meth-nn}

The deep model composes $L$ hidden layers,
\begin{equation}
\label{eq:mlp}
z^{(0)} = \phi(x), \qquad
z^{(\ell)} = \sigma\!\left(W^{(\ell)} z^{(\ell-1)} + b^{(\ell)}\right),
\qquad
p_\theta(\cdot \mid x) =
\operatorname{softmax}\!\left(W^{(L+1)} z^{(L)} + b^{(L+1)}\right),
\end{equation}
with rectified-linear activations $\sigma(u)=\max(u,0)$ and the
input map $\phi$ of Section~\ref{sec:meth-features}.  Depth is
searched over $L\in\{1,3,5,7\}$ anchored at the five-layer geometry
preferred by \citet{sirignano2021} (200 units, then 140 per layer);
layer widths are inherited from that paper rather than searched ---
depth and regularisation are the first-order capacity knobs --- with
a half/paper/double width sweep at the selected configuration as a
completeness check (Appendix~\ref{app:supplementary}).  Each hidden
unit is a learned ``soft feature''; depth lets the network build
feature hierarchies, which is where interaction effects such as
incentive$\times$credit-score live --- the model discovers them
rather than being told.

Three regularisation devices mirror the paper: \emph{dropout}
\citep{dropout} on hidden layers (rates $\{0, 0.2, 0.5\}$), an
$\ell_2$ penalty ($\lambda \in \{0, 10^{-5}, 10^{-4}\}$), and
\emph{early stopping} on the window's validation likelihood --- the
only use the validation year is put to.  Finally, the headline model
is an \emph{ensemble}: $K{=}8$ independently initialised networks
trained on resampled shard subsets, with predicted probabilities
averaged, $\bar p(j\mid x) = \tfrac1K \sum_k p_{\theta_k}(j \mid x)$.
Averaging reduces the variance component of error from stochastic
training and finite data; the marginal value of each additional
member is traced in an ensemble-size curve.  Because eight fits per
window is the dominant compute cost, the ensemble runs on five
regime-spanning key windows ($k\in\{2015, 2019, 2020, 2023, 2025\}$),
with the single best network everywhere.

\section{Evaluation}
\label{sec:meth-eval}

The headline criterion is the out-of-sample mean negative
log-likelihood,
\begin{equation}
\label{eq:nll}
\mathrm{NLL} \;=\; -\frac{1}{N_{\mathrm{test}}}
\sum_{i\in\mathrm{test}} \log p_\theta\!\left(y_i \mid x_i\right),
\end{equation}
on identical, frozen, unthinned test rows for every model.  It is the
right primary metric for two reasons: it evaluates the full predicted
distribution over seven destinations rather than a binarised event,
and it is a strictly proper scoring rule --- uniquely minimised in
expectation by the true conditional probabilities, so a model cannot
improve it by distorting its probabilities.  Two complementary lenses
guard against failure modes a single number hides:
\emph{discrimination}, measured for each origin--destination pair by
the one-versus-rest area under the ROC curve (the probability that a
randomly chosen loan-month that truly makes the transition receives a
higher predicted probability than one that does not); and
\emph{calibration}, measured by reliability diagrams comparing
predicted probabilities to realised frequencies within prediction
buckets.  A model can rank well yet be systematically over- or
under-confident; calibration is what the pool-level counting of
Section~\ref{sec:meth-pool} actually consumes.

\section{From loans to pools}
\label{sec:meth-pool}

A mortgage pool's cashflow is the sum of its loans' events; a
mortgage-backed security (MBS) is priced off exactly such sums.
\citet{sirignano2021} argue the nonlinearity advantage
\emph{compounds} rather than washes out at the pool level, because
linear misspecification correlates across loans sharing
characteristics.

\subsection{Multi-month prediction by matrix composition}
\label{subsec:meth-rollforward}

The monthly model is rolled forward by composing transition
matrices.  For loan $i$ alive at anchor month $t_0$ with state
distribution $\rho_0$ (a one-hot vector), the model supplies for each
future month $h$ a $7\times 7$ matrix $P^{(i)}_h$ whose four
non-terminal rows are the model's predictions at the
deterministically evolved covariates and whose terminal rows are the
identity (absorbing states).  Then
\begin{equation}
\label{eq:rollforward}
\rho_h \;=\; \rho_0 \prod_{h'=1}^{h} P^{(i)}_{h'},
\qquad
\mathbb{P}\!\left(\text{loan } i \text{ prepaid by } t_0{+}H\right)
= \left[\rho_H\right]_{\mathrm{Prepaid}},
\end{equation}
with horizon $H=12$ months.  Because terminal states are absorbing
and the full state distribution is tracked,
\eqref{eq:rollforward} is exact under the model --- no Monte Carlo
simulation is required.  Covariates evolve deterministically (age and
remaining term advance, seasonality follows the calendar) while
balances and the macroeconomic block are frozen at $t_0$: predictions
are explicitly \emph{conditional on the anchor-date environment}.
For ``ever reaches 60+ days past due within the horizon'' a
first-passage construction is used: delinquency states can be entered
and cured out of, so month-$H$ occupancy undercounts loans that
visited and recovered; making the target set absorbing in a modified
chain yields the probability of ever hitting it.

\subsection{Pools and pool-level prediction}
\label{subsec:meth-pools}

Pools are formed two ways at each window anchor: characteristic
buckets (credit score $\times$ original rate $\times$
loan-to-value cells, the criteria a securitisation structurer would
use) and random pools ($B$ pools of $N$ loans), which give a
bucketing-free comparison.  With per-loan horizon probabilities $q_i$
from \eqref{eq:rollforward}, the predicted event count in pool
$\mathcal{P}$ and its approximate distribution are
\begin{equation}
\label{eq:poolcount}
\widehat{C}_{\mathcal{P}} = \sum_{i\in\mathcal{P}} q_i,
\qquad
C_{\mathcal{P}} \;\dot\sim\;
\mathcal{N}\!\Bigl(\textstyle\sum_i q_i,\; \sum_i q_i(1-q_i)\Bigr),
\end{equation}
the Gaussian being the standard approximation to the
Poisson--binomial distribution of a sum of independent non-identical
Bernoulli variables (conditional independence given covariates is the
maintained assumption, inherited from the paper).  Models are
compared on predicted versus realised counts, with interval coverage
testing the distributional claim in \eqref{eq:poolcount}, not just
its mean.

\subsection{Economic translation: prepayment speed, cashflow timing, and price}
\label{subsec:meth-econ}

The final step converts counts into valuation units; each term is
defined as it first appears.  Averaging the loans' monthly prepayment
probabilities with each loan weighted by its unpaid principal balance
gives the pool's \textbf{prepayment-speed path}: the model's
month-by-month forecast of what fraction of the pool's remaining
principal prepays --- the units and weighting that pool cashflows
actually obey, since a large loan returns proportionally more
principal than a small one.  (Market shorthand, used hereafter: the
monthly fraction is the \emph{single monthly mortality}, SMM;
annualised it is the \emph{conditional prepayment rate}, CPR.)  The
realised path is computed identically from the panel.  A
deterministic cashflow engine turns any speed path into pool
cashflows, treating the pool as a \emph{pass-through} (investors
receive the pool's interest and principal as it is paid): scheduled
amortisation plus forecast prepayments, month by month, with the
final month's speed held constant beyond the horizon.  From the
cashflows follow two valuation summaries: \textbf{weighted-average
life} --- the average time a dollar of principal remains outstanding,
i.e.\ the timing of money returned --- and \textbf{price}, the
discounted value of the cashflows, with every model's cashflows
discounted at the same fixed rate so that any price difference comes
from the prepayment forecasts alone.  The comparison metrics are then
a speed error (percentage points of CPR), a timing error (months of
weighted-average life), and a price error (per 100 of face value),
per pool, model, and anchor.  The engine is unit-tested against two
closed forms (with no prepayment a pool is a textbook annuity; a
constant speed has a closed-form survival schedule).  With the
macroeconomic environment frozen at the anchor and a common fixed
discount rate, the exercise isolates the \emph{prepayment-model
component} of valuation error --- deliberately simpler than the
market's full option-adjusted-spread machinery.

\paragraph{Why pricing rather than portfolio sorts.}
\citet{sirignano2021} demonstrate economic significance with a
portfolio-sort exercise: rank pools by predicted prepayment risk,
form best- and worst-decile portfolios, and compare their realised
outcomes.  This study runs the pricing exercise instead, and the
choice deserves a plain explanation.  The two designs test different
properties of the same predictions.  A sort uses only the
\emph{ordering} of pools --- it is blind to probability levels, so a
model whose hazards are uniformly twice the truth sorts exactly as
well as a perfectly calibrated one.  Pricing uses the \emph{levels}:
weighted-average life and price depend on the actual magnitude of the
predicted prepayment path.  Ranking-only evidence is the natural
choice where levels are inherently untrustworthy (the classic
asset-pricing setting); here, by contrast, the transition
probabilities are explicitly calibration-checked
(Section~\ref{sec:meth-eval}), so a ranking-only exercise would
discard the model's strongest property --- and the practical consumer
of a pool-level mortgage model (pricing a pool, marking it to market,
hedging a position in it) consumes levels, not ranks.  Running both
would add little: each is a deterministic re-presentation of the same
per-loan probabilities, their respective properties are already
tested more rigorously upstream (discrimination by the AUC tables,
calibration by the reliability diagrams), and the bucket-level
price-error map conveys what a sort would dramatise --- \emph{where}
the linear model goes wrong --- while also conveying by \emph{how
much}, in the units the market trades.  The sort exercise is
therefore noted as future work rather than duplicated.

\section{Reproducibility}
\label{sec:meth-repro}

Every model fit writes a self-describing run record (configuration,
code version, data-manifest hash, window identity, metrics), so any
number reported in Chapter~\ref{chap:results} can be traced to the
exact code and data slice that produced it.  Correctness-critical
components carry unit tests: the leakage mask, the importance
weights, the mark-to-market LTV construction, the roll-forward engine
against a closed-form chain, and the cashflow engine against the
annuity closed forms.  Appendix~\ref{app:reproducibility} records the
computing environment and data snapshot conventions.
